\section{性能分析}
\begin{code}
\begin{verbatim}
import torch
import torchvision.models as models
from torch.profiler import profile, record_function, ProfilerActivity

model = models.resnet18().cuda()
inputs = torch.randn(5, 3, 224, 224, device="cuda")

with profile(
    activities=[ProfilerActivity.CPU, ProfilerActivity.CUDA],
    profile_memory=True,
) as prof:
    with record_function("model_inference"):
        model(inputs)

prof.export_chrome_trace("trace.json")
print(prof.key_averages().table(sort_by="cuda_time_total", row_limit=10))
\end{verbatim}
\end{code}

\href{https://ui.perfetto.dev}{perfetto trace viewer}，如果需要比Pytorch Profiler更底层的GPU性能分析，
可以考虑使用 Nsight Compute。如果需要CPU性能分析，可以尝试Py-Spy和strace。for-loop，for-each和fuse

设置参数，拒绝PyTorch分配大显存给小参数。
\begin{code}
\begin{verbatim}
export PYTORCH_CUDA_ALLOC_CONF=max_split_size_mb:128
\end{verbatim}
\end{code}

\begin{code}
\begin{verbatim}
import torch

torch.cuda.memory._record_memory_history()

with torch.inference_mode():
    shape = [156, 1024, 1024, 1]
    x1 = torch.randn(shape, device="cuda:0")
    x2 = torch.randn(shape, device="cuda:0")
    y = x1 * x2

torch.cuda.memory._dump_snapshot("vram_profile.pickle")
\end{verbatim}
\end{code}
\href{https://pytorch.org/memory_viz}{PyTorch Memory Visualization}